% Generated draft: requires mathematical and source review.
\documentclass[11pt]{article}
\usepackage[margin=25mm]{geometry}
\usepackage{fontspec}
\setmainfont{Latin Modern Roman}
\usepackage{amsmath,amssymb,mathtools,unicode-math}
\setmathfont{Latin Modern Math}
\usepackage{enumitem,needspace,xurl,hyperref}
\setlength{\parindent}{0pt}
\setlength{\parskip}{5pt}
\setlength{\emergencystretch}{3em}
\newcommand{\N}{\mathbb{N}}
\newcommand{\Z}{\mathbb{Z}}
\newcommand{\Q}{\mathbb{Q}}
\newcommand{\R}{\mathbb{R}}
\newcommand{\C}{\mathbb{C}}
\newcommand{\F}{\mathbb{F}}
\newcommand{\A}{\mathbb{A}}
\newcommand{\PP}{\mathbb{P}}
\newcommand{\E}{\mathbb{E}}
\newcommand{\eps}{\varepsilon}
\newcommand{\dd}{\,\mathrm{d}}
\begin{document}
\section*{UnsolvedMath ID: 30003617}
\textbf{Source problem code:} catalogue entry
\subsection{Definitions and mathematical statement}
A topological dynamical system consists of a topological group \(G\), a nonempty topological space \(X\), and a continuous action \(G\times X\to X\), written \((g,x)\mapsto gx\). It is minimal if every orbit is dense:
\[
\forall x\in X,\qquad \overline{\{gx:g\in G\}}=X.
\]
The group is commutative if \(gh=hg\) for all \(g,h\in G\).

To explain one standard meaning of tameness, suppose additionally that \(X\) is compact and metrizable. For \(g\in G\), let \(T_g(x)=gx\), and define the enveloping semigroup by
\[
E(X,G)=\overline{\{T_g:g\in G\}}^{\,X^X},
\]
where \(X^X\) carries the topology of pointwise convergence. In this setting, tameness can be expressed by requiring every element of \(E(X,G)\) to be of Baire class one, meaning a pointwise limit of a sequence of continuous maps \(X\to X\). This explanation is conditional: the supplied record does not specify compactness, metrizability, the topology of \(G\), or its intended definition of tameness, so these additional hypotheses are not adopted as restrictions on the source question.

For a discrete group \(G\), amenability means that there exists a positive linear functional \(m:\ell^\infty(G;\R)\to\R\) such that \(m(\mathbf 1)=1\) and
\[
m(f\circ L_g)=m(f)\qquad(g\in G,\ f\in\ell^\infty(G;\R)),
\qquad L_g(h)=gh.
\]
The source does not specify whether this discrete-group convention or a topological-group notion of amenability is intended.

\textbf{Source question.} What happens for tame minimal dynamical systems when the acting group is not assumed to be commutative or even amenable?

This asks for an investigation beyond commutative acting groups, including the possibility of nonamenable acting groups. It does not require every acting group under consideration to be noncommutative or nonamenable. The supplied record identifies no particular structural conclusion, classification, or existence assertion to establish. Consequently, a uniquely specified theorem or conjecture with complete hypotheses and conclusion cannot be extracted from this record alone; the question remains a research direction with the unresolved qualifications stated above.
\subsection{Short English statement}
The question asks how the theory of tame minimal dynamical systems extends when commutativity and amenability of the acting group are not assumed. The supplied record leaves both the precise class of systems and the desired conclusion unspecified, so further source clarification is needed for a fully precise conjecture.
\subsection{Sources}
\begin{itemize}
\item Supplied local UnsolvedMath registry record, permanent ID \(30003617\), source problem code \texttt{catalogue entry}, title: \emph{Tame Dynamics of Nonamenable Group Actions}. The record belongs to the Oberwolfach Reports Open Problems collection and gives the proposed year as 2017. No individual proposer is supplied.
\item Supplied report citation: \emph{Spectral Structures and Topological Methods in Mathematical Quasicrystals} (2018), Oberwolfach Reports, pp.~2781--2845, DOI \texttt{10.4171/owr/2017/46}. Supplied source and literature-reference URL: \url{https://doi.org/10.4171/owr/2017/46}. The literature-reference label is ``Oberwolfach Report 46 2017.'' The background describes these reports as documenting mathematical workshops and their open-problem discussions.
\item Literature assessment dated exactly \texttt{2026-08-21}: the supplied curation labels the problem open and says that the extracted question concerns extension to noncommutative or nonamenable acting groups, with no general theory or resolution supplied. This is supplied dated curation, not an independently verified literature assessment. The background explicitly requests independent verification of time-sensitive status and cited claims.
\item The registry labels the statement \texttt{corrected\_verified} and reports a manual semantic, typographic, and LaTeX review claiming self-containedness and preservation of the mathematical target. These are supplied editorial assertions; they do not resolve the missing hypotheses or identify a precise mathematical conclusion.
\end{itemize}
\end{document}
